% Part: sets-functions-relations
% Chapter: infinite
% Section: dedekind
%
\documentclass[../../../include/open-logic-section]{subfiles}

\begin{document}

\olfileid{sfr}{infinite}{dedekind}
\olsection{डेडेकिंड बीजसंरचना}

नैसर्गिक संख्या केवळ अनंत असाव्यात असेच आपल्याला वाटत नाही; त्यांना काही
(बीजगणितीय) गुणधर्म असावेत असेही आपल्याला वाटते: बेरीज, गुणाकार इत्यादी
क्रियांखाली त्यांनी योग्य रीतीने वागले पाहिजे.

\emph{उत्तरवर्ती फलनाची} कल्पना मूलभूत मानून, पुढील गुणधर्म असलेल्या
वस्तू म्हणजे संख्या असे त्यांचे वैशिष्ट्य सांगावे, ही डेडेकिंड यांची
कल्पना होती:
\begin{enumerate}
	\item अशी $0$ ही संख्या आहे जी कोणत्याही संख्येची उत्तरवर्ती नाही
	\\म्हणजे, $0 \notin \ran{s}$
	\\म्हणजे, $\forall x\ s(x) \neq 0$
	\item भिन्न संख्यांच्या उत्तरवर्ती संख्या भिन्न असतात
	\\म्हणजे, $s$ हे !!a{injection} आहे
	\\म्हणजे, $\forall x \forall y (s(x) = s(y) \lif x = y)$
	\item\ollabel{repeatedapplication} प्रत्येक संख्या $0$ पासून
	उत्तरवर्ती फलन वारंवार लावून मिळते.
\end{enumerate}
पहिल्या दोन अटी प्रथम-क्रम तर्कशास्त्र वापरून सहज हाताळता येतात (वर
पाहा). पण केवळ प्रथम-क्रम तर्कशास्त्र वापरून \olref{repeatedapplication}
ही अट हाताळता येत नाही. \olref{repeatedapplication} ही अट पुढीलप्रमाणे
संचसैद्धान्तिक रूपात पुन्हा मांडणे ही डेडेकिंड यांची निर्णायक कामगिरी होती:
\begin{enumerate}
	\item[3$'$.] नैसर्गिक संख्यांचा संच हा \emph{उत्तरवर्ती फलनाखाली
	बंद} असलेला लघुतम संच आहे: म्हणजे, संचातील कोणत्याही !!{element} वर
	$s$ लावल्यास संचातीलच दुसरा !!{element} मिळतो.
\end{enumerate}
पण हे आपल्याला सावकाश आणि स्पष्टपणे मांडावे लागेल.

\begin{defn}\ollabel{Closure}
	कोणत्याही फलन $f$ साठी, संच $X$ हा $f$-\emph{बंद} असतो जर आणि तरच
	$(\forall x \in X)f(x) \in X$. आता कोणत्याही $o$ साठी अशी व्याख्या करू:
	$$\closureofunder{f}{o} = \bigcap\Setabs{X}{o \in X\text{ आणि }X
\text{ हा $f$-बंद आहे}}$$
\end{defn}

म्हणून $\closureofunder{f}{o}$ हा $o$ हा !!a{element} असलेल्या सर्व
$f$-बंद संचांचा छेद आहे. त्यामुळे सहजज्ञानानुसार,
$\closureofunder{f}{o}$ हा $o$ हा !!a{element} असलेला \emph{लघुतम}
$f$-बंद संच आहे. पुढील निष्कर्ष हे सहजज्ञान नेमकेपणाने मांडतो:
\begin{lem}\ollabel{closureproperties}
	कोणतेही फलन $f$ आणि कोणतेही $o \in A$ यांसाठी:
	\begin{enumerate}
		\item\ollabel{closurehaselem} $o \in \closureofunder{f}{o}$; आणि
		\item\ollabel{closureclosed} $\closureofunder{f}{o}$ हा $f$-बंद आहे; आणि
		\item\ollabel{closuresmallest} $X$ हा $f$-बंद असून $o \in
		X$ असेल, तर $\closureofunder{f}{o} \subseteq X$
	\end{enumerate}
\end{lem}

\begin{proof}
$o$ हा !!a{element} असलेला किमान एक $f$-बंद संच आहे, हे लक्षात घ्या;
तो म्हणजे $\ran{f}\cup \{o\}$. म्हणून अशा \emph{सर्व} संचांचा छेद
$\closureofunder{f}{o}$ अस्तित्वात आहे. आता आपल्याला
\olref{closurehaselem}--\olref{closuresmallest} तपासावे लागतील.

\olref{closurehaselem} बाबत: छेदातील सर्व संचांत $o$ हा !!a{element}
असल्यामुळे $o \in \closureofunder{f}{o}$.

\olref{closureclosed} बाबत: $x \in \closureofunder{f}{o}$ असे समजा.
म्हणून $o \in X$ आणि $X$ हा $f$-बंद असेल, तर $x \in X$; आणि $X$ हा
$f$-बंद असल्यामुळे $f(x) \in X$. म्हणून
$f(x) \in \closureofunder{f}{o}$.

\olref{closuresmallest} बाबत: सर्वसाधारणपणे, $X \in C$ असेल, तर
$\bigcap C \subseteq X$.
\end{proof}

याचा उपयोग करून आपण असे म्हणू शकतो:

\begin{defn}
संच $A$, फलन $f \colon A \to A$ आणि असा काही $o \in A$ यांनी मिळून
\emph{डेडेकिंड बीजसंरचना} बनते, जर:
	\begin{enumerate}
		\item \ollabel{ded:proper} $o \notin \ran{f}$
		\item \ollabel{ded:injection} $f$ हे !!a{injection} आहे
		\item \ollabel{ded:closure} $A = \closureofunder{f}{o}$
	\end{enumerate}
\end{defn}

$A = \closureofunder{f}{o}$ असल्यामुळे, आधीच्या निष्कर्षावरून $A$ हा
$o$ हा !!a{element} असलेला लघुतम $f$-बंद संच आहे. डेडेकिंड बीजसंरचना
स्पष्टपणे डेडेकिंड-अनंत असते; व्याख्येतील \olref{ded:proper} आणि
\olref{ded:injection} ही कलमे पाहा. पण कोणत्याही डेडेकिंड-अनंत संचाचे
डेडेकिंड बीजसंरचनेत रूपांतर करता येते, ही अधिक रोचक गोष्ट आहे.

\begin{thm}\ollabel{thm:DedekindInfiniteAlgebra}
डेडेकिंड-अनंत संच अस्तित्वात असेल, तर डेडेकिंड बीजसंरचना अस्तित्वात असते.
\end{thm}

\begin{proof}
$D$ हा डेडेकिंड-अनंत असू द्या. मग एकास-एक फलन $g \colon D \to D$
आणि घटक $o  \in D \setminus \ran{g}$ अस्तित्वात आहेत. आता $A =
\closureofunder{g}{o}$ असे मानू; \olref{closureproperties} नुसार $A$
अस्तित्वात आहे आणि $o \in A$. $f = \funrestrictionto{g}{A}$ असे मानू.
$A, f, o$ हे मिळून डेडेकिंड बीजसंरचना बनतात, हे आपण दाखवू.

\olref{ded:proper} बाबत: $o \notin \ran{g}$ आणि $\ran{f}
\subseteq \ran{g}$, म्हणून $o\notin \ran{f}$.

\olref{ded:injection} बाबत: $g$ हे $D$ वरील एकास-एक फलन आहे; म्हणून
$f \subseteq g$ हेही एकास-एक फलन असले पाहिजे.

\olref{ded:closure} बाबत: \olref{closureproperties} नुसार $A$ हा
$g$-बंद आहे; त्यामुळे विशेषतः $A$ हा $f$-बंद आहे. म्हणून
\olref{closureproperties} नुसार $\closureofunder{f}{o} \subseteq A$.
तसेच $\closureofunder{f}{o}$ हा $f$-बंद आहे आणि
$f = \funrestrictionto{g}{A}$ असल्यामुळे $\closureofunder{f}{o}$ हा
$g$-बंद आहे. म्हणून \olref{closureproperties} नुसार
$A \subseteq \closureofunder{f}{o}$.
\end{proof}

\end{document}
